% ===========================================================================
\section{A Bayesian Interpretation of the Population Model}
\label{sec:bayesian}

The population model of Section~\ref{sec:equilibria} reads
$\mathbf{f}$ as the frequency of strategies in an evolving population
and Equation~\eqref{eq:softmax} as selection pressure.
The motivating scenario suggests a different reading.
The authors of strategies are technologically advanced actors ---
civilizations, or organizations with access to strong AI ---
and a strategy is the endpoint of a deliberate investigation,
not a variant surviving rounds of selection.
This section shows that the same fixed-point machinery arises from
Bayesian reasoning by such authors, with two benefits:
the inverse temperature $\beta$ acquires a rational interpretation,
and the model extends naturally to authors who differ in how much
they value their own payoff versus the total welfare of the
tournament.

\subsection{From selection to belief}

An entrant choosing a strategy cannot distinguish the other entrants
before observing their submissions: the field is \emph{exchangeable}.
By de Finetti's theorem \citep{definetti1937prevision}, exchangeability
is equivalent to positing a latent distribution $\mathbf{f}$ from which
the field is drawn independently.
Reasoning about opponents is therefore reasoning about $\mathbf{f}$ ---
the frequency vector of Section~\ref{sec:equilibria}, reinterpreted as
a \emph{belief} about the distribution of strategies an entrant will face.

The belief is self-referential.
Each author expects the field to be weighted toward strategies that
perform well against the field itself, because every other author
reasons the same way.
An equilibrium belief is one that reproduces itself:
the distribution of strategies that optimizing authors choose,
given the belief, is the belief.
This is a rational-expectations condition in the sense of
Bayesian games with a common prior \citep{harsanyi1967games},
and it can equally be read as the limit of iterated deliberation:
posit a set of plausible strategies, evaluate how they interact,
shift credence toward the ones that perform well, and repeat.

\subsection{The candidate pool}

In the motivating scenario a strategy is neither a formal program nor a
natural-language document; it is a disposition arrived at by
investigation, of which the artifacts in our tournaments are proxies.
We therefore do not attempt a prior over any unrestricted space of
strategy artifacts.
Instead we work with an empirical \emph{candidate pool}
$\mathcal{P}$: the strategies submitted to our tournaments together
with strategies elicited from AI models
(Section~\ref{sec:archetypes}), and we place prior distributions over
$\mathcal{P}$; each tournament's strategy pool is a subset
$\mathcal{S} \subseteq \mathcal{P}$.
The analytical questions then become robustness questions:
whether and how the equilibrium answers change as the prior over
$\mathcal{P}$ shifts, and as $\mathcal{P}$ itself grows.

Working relative to a pool is not merely a practical concession.
No strategy is even weakly dominant over the unrestricted space of
possibilities, because for any strategy $s$ there exists a strategy
that recognizes $s$ exactly and conditions its decision on that
recognition.
Dominance results (Section~\ref{sec:dominance}) are therefore
necessarily statements about a candidate pool $\mathcal{P}$ and a
prior over it, never about the game in the abstract.
The same holds for equilibrium multiplicity: when the fixed-point
equation has several solutions, which one is reached is determined by
the initial credence --- the prior --- and the multi-start procedure of
Section~\ref{sec:equilibria} is, under this reading, a scan over
priors.

\subsection{Author types}
\label{sec:author-types}

Perhaps the most consequential difference between authors is how much
they value their own payoff against the total welfare of the
tournament.
We index types by $\lambda \in [0, 1]$.
Alongside the payoff matrix $A$ of Equation~\eqref{eq:payoff}, define
the \emph{welfare matrix} $W$ by
\begin{equation}
  W_{ij} = u_d(d_{ij}) + u_r(d_{ij}) + u_d(d_{ji}) + u_r(d_{ji}),
  \label{eq:welfare-matrix}
\end{equation}
the total welfare generated by both rounds between strategies $i$ and
$j$, counting both parties.
A type-$\lambda$ author evaluates strategy $i$ against belief
$\mathbf{f}$ by
\begin{equation}
  U^\lambda_i(\mathbf{f})
  = (1 - \lambda)\, (A\mathbf{f})_i + \lambda\, (W\mathbf{f})_i .
  \label{eq:type-utility}
\end{equation}
Type $\lambda = 0$ maximizes own payoff, recovering
Equation~\eqref{eq:reward}; type $\lambda = 1$ maximizes total
welfare.
(Total welfare includes rounds not involving the chosen strategy, but
for an individual author those rounds are unaffected by the choice and
drop out of the comparison; they return in
Section~\ref{sec:bayesian-bloc}.)

The direct trade-off between the objectives has a closed form in the
constants of Section~\ref{sec:properties}: sharing with a recipient
who will not reciprocate costs $(1-\lambda)\,\delta$ of weighted own
payoff and contributes $\lambda\, g$ of weighted welfare, so a type
prefers unconditional generosity toward non-reciprocators only if
\begin{equation}
  \lambda > \lambda^* = \frac{\delta}{g} = \frac{1}{\rho},
  \label{eq:lambda-threshold}
\end{equation}
which for $E = 60$ gives $\lambda^* \approx 0.197$.
Below the threshold, all types favor conditional cooperation on direct
payoffs alone; the equilibrium effects computed below sharpen this
further.

\paragraph{What is the most self-interested strategy?}
The type index turns loose questions about motive into precise ones:
which strategy should a type-$\lambda$ author submit?
For the own-payoff author ($\lambda = 0$) the preceding sections
already contain the answer, and it is not the plain reading.
\textbf{Selfish} --- take from everyone, maximizing each round in
isolation --- is weakly dominated by \textbf{Cooperation coalition}
(Section~\ref{sec:dominance-results}): whatever the field,
indiscriminate taking earns no more than selective taking, and
against some fields strictly less.
The surviving candidates are conditional: \textbf{Mirror} tops the
leaderboard, and the own-payoff equilibrium concentrates on the
enforcement core (Section~\ref{sec:results-softmax}).
Self-interest, pursued rationally against other rational
participants, chooses conditional cooperation --- reserving
\textup{\textsc{take}} for defectors and their subsidizers.

\paragraph{What is the most altruistic strategy?}
For the $\lambda = 1$ author the answer is less settled, and the
question is worth posing carefully.
Four candidates come with intuitive cases:
(a)~\textbf{Generous}: unconditional giving is the plain reading of
altruism;
(b)~a conditional cooperator such as \textbf{Mirror}: by making
defection unprofitable it steers rational participants away from
\textbf{Selfish} --- altruism through incentives rather than gifts;
(c)~\textbf{Cooperation coalition}: it pushes the field away from
\textbf{Selfish} \emph{and} away from \textbf{Generous}, and the
second push reinforces the first, since every unconditional
cooperator it drives out is a subsidizer whose absence lowers a
defector's revenue (Proposition~\ref{prop:drift});
(d)~\textbf{Chivalry}: conditional cooperation \emph{plus}
protection --- it deters defection while keeping the innocents it
defends in the field.
Note that (c) and (d) embody the two antagonistic second-order norms
of Section~\ref{sec:second-order}: they disagree precisely over
whether (a) is to be harvested or defended.
None of this can be settled by reading the strategy texts, because
each case is a claim about how rational co-players \emph{respond} in
equilibrium.
The bloc formalism of Section~\ref{sec:bayesian-bloc} turns the
question into a computation, and the answer is that the plain
reading finishes near the bottom.

\paragraph{Neither objective is served by causal maximization.}
Both answers defy the decision-level logic of their own objectives.
Fix any encounter and evaluate the
\textup{\textsc{share}}/\textup{\textsc{take}} decision by causal
lights, holding every other decision fixed:
\textup{\textsc{take}} raises the dictator's own payoff by $\delta$
in every case, and \textup{\textsc{share}} raises the round's
welfare by $g - \delta > 0$ in every case.
A causal maximizer of own payoff therefore always takes, and a
causal maximizer of total welfare always shares: applied decision by
decision, causal reasoning recommends exactly the two unconditional
strategies --- the plain readings rejected above, both weakly
dominated (Section~\ref{sec:dominance-results}).
The strategies that actually optimize each terminal goal commit to
decisions that are causally suboptimal for that very goal:
\textbf{Mirror} and the coalition share where taking would earn
$\delta$ more, and the welfare bloc's best choices take from
violators, burning $g - \delta$ of the very surplus their objective
counts, to sustain deterrence.
What licenses these commitments is that the OSDG offers no decision
node at which to deviate unobserved: the decision is computed by the
oracle from the artifact, so a strategy that would take where
\textbf{Mirror} shares is a \emph{different artifact}, and every
conditional opponent treats it differently.
The correlation between one's disposition and others' decisions is
evaluative rather than causal --- and it is where all the payoff
lives.
This is the sense in which the open-strategy setting rewards
functional rather than causal decision theory
(Section~\ref{sec:related}).

\subsection{Random utility and the mixture fixed point}

Even ideally rational authors of the same type $\lambda$ differ in
idiosyncratic ways --- private context, secondary objectives,
judgment calls that the type index does not capture.
Following the random-utility model of discrete choice
\citep{mcfadden1974conditional}, we add to
$U^\lambda_i$ an author-specific shock with Gumbel tails, under which
the distribution of choices made by type-$\lambda$ authors is the
logit
\begin{equation}
  q^\lambda_i(\mathbf{f}) \propto
  e^{\beta\, U^\lambda_i(\mathbf{f})}.
  \label{eq:logit-choice}
\end{equation}
Here $\beta$ is the precision of the common, payoff-driven component
of the objective relative to the idiosyncratic component --- not a
bound on rationality and not a strength of selection.
Advanced authors correspond to large $\beta$, but $\beta$ remains
finite as long as authors differ at all.

With type masses $m_t$ summing to one, the rational-expectations
equilibrium is the fixed point
\begin{equation}
  \mathbf{f}^* = \sum_t m_t\,
  \mathrm{softmax}\!\big(\beta\, U^{\lambda_t}(\mathbf{f}^*)\big).
  \label{eq:mixture-fixed-point}
\end{equation}
A single type with $\lambda = 0$ recovers
Equation~\eqref{eq:softmax} exactly: the softmax equilibrium of
Section~\ref{sec:equilibria} is the special case of homogeneous
own-payoff authors, and is formally a logit quantal-response
equilibrium \citep{mckelvey1995quantal} of the submission game.
Existence follows as before: the right-hand side of
Equation~\eqref{eq:mixture-fixed-point} is a continuous self-map of
the simplex, so Brouwer's theorem applies, and every computation of
Section~\ref{sec:equilibria} carries over with $U^\lambda$ in place
of $R$.

\subsection{Illustration: shifting the prior over the pool}
\label{sec:bayesian-illustration}

Applying the model to the Tournament~9 pool
illustrates how the answers move --- and fail to move --- as the prior shifts.

\paragraph{Objectives pool at equilibrium.}
With author types in proportions $0.4$ own-payoff, $0.4$ mixed
($\lambda = 0.5$), $0.2$ total-welfare, all three types concentrate on
the same cooperative core
(\textbf{Mirror}, \textbf{Conditional cooperator},
\textbf{Cooperation coalition}), diverging only marginally; the
bottom row of Figure~\ref{fig:trajectories} traces the convergence of
the mixture fixed point, Equation~\eqref{eq:mixture-fixed-point}, at
$\beta = 5$ and $\beta = 20$.
Because every mutual \textup{\textsc{share}} creates a surplus
$g - \delta > 0$ that pays both parties, strategies maximizing mutual
cooperation are near-optimal for every $\lambda$; objective
heterogeneity is nearly unidentifiable from submissions.

\paragraph{Prior shifts select among equilibria.}
At large $\beta$ the own-payoff fixed point is not unique: the two
basins of Section~\ref{sec:results-basins} realize the enforcement
and protection configurations of the antagonistic second-order norms
(Section~\ref{sec:second-order}).
Both configurations are fully cooperative on their support --- every
round a mutual \textup{\textsc{share}} --- and both attain the
first-best welfare $V = 2\ln(1 + E/2)$ in the sharp limit.
At $\beta = 20$ the tie is not exact: the protection fixed point
retains a remnant of its rival norm (\textbf{Cooperation coalition},
mass $0.008$), and the remnant's feud with \textbf{Chivalry} and
harvesting of the subsidizers cost $0.3\%$ of welfare --- a gap that
vanishes as $\beta$ grows.
Welfare therefore barely separates the basins; which one obtains is
decided by the prior over the pool, and Figure~\ref{fig:basin-map}
is, under this reading, a map of that decision: the initial credence
assigned to the two rival norms determines the equilibrium that
iterated deliberation reaches.
Under the type mixture above, by contrast, multi-start sampling
($300$ Dirichlet starts) finds a single fixed point at both
temperatures: at these type masses the protection configuration is no
longer an attractor.
Growing the pool has the same character: adding an AI-elicited
strategy can create or destroy dominance relations and shift basin
boundaries, so we report equilibrium conclusions together with the
pool and prior against which they were computed.
